% Einheit 3 von 4 · Sachaufgaben (Katalog Einheit 4)
\hypertarget{e3}{}\einheitenkopf{Einheit 3 von 4 · Sachaufgaben}
\zweigzeile{Hier lernst du, zu einer Sachaufgabe eine quadratische Gleichung aufzustellen, sie zu lösen und die passende Lösung auszuwählen · neu in diesem Jahr, je nach Buch erst Klasse 10 · keine P10-Aufgabe · baut auf: Wurzelziehen und Lösbarkeit (Einheit 1), Normalform und p-q-Formel (Einheit 2), Gleichungen aufstellen}
\renewcommand{\theaufgabe}{\arabic{aufgabe}}\setcounter{aufgabe}{21}

\begin{aufgabe}{Ich erkenne, welche Lösung zur Frage passt. Zu jeder Frage stehen die beiden Lösungen der Gleichung da. Kreuze an, welche passt. Du musst nichts rechnen.}
\begin{geruest}
\gz{a}{Seitenlänge eines Beets in m: \ $x_1 = 7$, \ $x_2 = -3$}{\kreuz{$x_1$}\kreuz{$x_2$}\kreuz{beide}}
\gz{b}{Anzahl der Gäste: \ $x_1 = -12$, \ $x_2 = 9$}{\kreuz{$x_1$}\kreuz{$x_2$}\kreuz{beide}}
\gz{c}{gesuchte Zahl im Zahlenrätsel: \ $x_1 = 4$, \ $x_2 = -4$}{\kreuz{$x_1$}\kreuz{$x_2$}\kreuz{beide}}
\gz{d}{Fallzeit eines Steins in s: \ $x_1 = 2{,}5$, \ $x_2 = -1{,}2$}{\kreuz{$x_1$}\kreuz{$x_2$}\kreuz{beide}}
\gz{e}{Temperatur in °C: \ $x_1 = -5$, \ $x_2 = 3$}{\kreuz{$x_1$}\kreuz{$x_2$}\kreuz{beide}}
\end{geruest}
\end{aufgabe}

\begin{aufgabe}{Ich kann ein Zahlenrätsel mit einer quadratischen Gleichung lösen. Die Gleichung steht dabei. Löse sie und gib alle Zahlen an, die passen.}
\beispiel{\text{Das Quadrat einer Zahl ist } 144{:} \quad x^2 &= 144 \\ x_1 &= 12, \quad x_2 = -12 \\ &\text{Die Zahl ist } 12 \text{ oder } {-12}.}
\begin{teile}
\teil Das Quadrat einer Zahl ist $225$. \quad $x^2 = 225$ \quad \feld{x_1}\feld{x_2}
\teil Das Quadrat einer Zahl ist $900$. \quad $x^2 = 900$ \quad \feld{x_1}\feld{x_2}
\teil Das Doppelte des Quadrats einer Zahl ist $50$. \quad $2x^2 = 50$ \quad \feld{x_1}\feld{x_2}
\teil Das Quadrat einer Zahl, vermehrt um $7$, ergibt $71$. \quad $x^2 + 7 = 71$ \quad \feld{x_1}\feld{x_2}
\end{teile}
Stelle die Gleichung selbst auf und löse sie.
\begin{teile}
\teil Subtrahiert man vom Quadrat einer Zahl $19$, erhält man $81$.\\ Gleichung: \leerfeld \feld{x_1}\feld{x_2}
\teil Das Produkt einer Zahl mit ihrem Nachfolger ist $72$. Gib beide Zahlenpaare an.\\ Gleichung: \leerfeld \feld{x_1}\feld{x_2}\\ Zahlenpaare: \leerfeld \leerfeld
\end{teile}
\end{aufgabe}

\begin{aufgabe}{Ich kann eine Rechteckaufgabe mit einer quadratischen Gleichung lösen. Im Beispiel ist ein Rechteck $3\,$m länger als breit und hat den Flächeninhalt $40\,\text{m}^2$.}
\beispiel{x \cdot (x + 3) &= 40 && \text{Breite } x, \text{ Länge } x + 3 \\ x^2 + 3x - 40 &= 0 \\ x_{1,2} &= -1{,}5 \pm \sqrt{2{,}25 + 40} = -1{,}5 \pm 6{,}5 \\ x_1 &= 5, \quad x_2 = -8 && {-8} \text{ ist keine Länge.} \\ &\text{Das Rechteck ist } 5\,\text{m breit und } 8\,\text{m lang.}}
Ein Rechteck ist $2\,$cm länger als breit und hat den Flächeninhalt $48\,\text{cm}^2$.

\rechteck[punkte=]{4}{2.4}

\begin{teile}
\teil Beschrifte die Seiten der Skizze mit $x$ und $x + 2$. Stelle die Gleichung auf.\\ Gleichung: \leerfeld \leerfeld
\teil Löse die Gleichung. \quad \feld{x_1}\feld{x_2}
\teil Welche Lösung passt zur Frage? Begründe.
\teil Schreibe einen Antwortsatz mit Einheit.
\teil Mache die Probe am Rechteck.
\end{teile}
\end{aufgabe}

\begin{aufgabe}{Ich kann eine Rechteckaufgabe mit einer quadratischen Gleichung lösen – weiter.}
\begin{teile}
\teil Ein quadratisches Foto bekommt ringsum einen $3\,$cm breiten Rahmen. Mit Rahmen ist das Bild $625\,\text{cm}^2$ groß. Wie lang ist eine Seite des Fotos?\\ Seitenlänge des Fotos: \leerfeld[cm]
\teil Ein rechteckiges Beet ist $5\,$m länger als breit. Man vergrößert beide Seiten um $1\,$m; dann ist das Beet $84\,\text{m}^2$ groß. Wie breit und wie lang war das Beet vorher? Stelle eine Gleichung auf, löse sie und gib die sinnvolle Lösung mit Einheit an.
\end{teile}
\end{aufgabe}

\begin{aufgabe}{Ich finde den Fehler in einer Rechteckaufgabe.}
\begin{teile}
\teil Ein Rechteck ist $4\,$m länger als breit und hat den Flächeninhalt $45\,\text{m}^2$. Emma rechnet so:
\rechnung{x \cdot (x + 4) &= 45 \\ x^2 + 4x - 45 &= 0 \\ x_1 &= 5, \quad x_2 = -9 \\ &\text{Die Breite ist } 5\,\text{m oder } {-9}\,\text{m.}}
Finde den Fehler, benenne ihn und korrigiere die Antwort.
\teil Ein Rechteck ist $5\,$m länger als breit und hat den Flächeninhalt $14\,\text{m}^2$. Berechne Breite und Länge.
\teil Ein Rechteck ist $3\,$cm länger als breit und hat den Flächeninhalt $54\,\text{cm}^2$. Paul stellt diese Gleichung auf:
\rechnung{2x + 2 \cdot (x + 3) &= 54}
Finde den Fehler, benenne ihn und korrigiere die Gleichung. Berechne dann Breite und Länge.
\teil Ein Rechteck ist $1\,$m länger als breit und hat den Flächeninhalt $30\,\text{m}^2$. Berechne Breite und Länge.
\end{teile}
\end{aufgabe}

\begin{aufgabe}{Ich kann begründen, wann beide Lösungen passen und wann nur eine. Zu beiden Fragen gehört die Gleichung $x^2 + 4x = 32$ mit den Lösungen $4$ und $-8$.}
\begin{teile}
\teil Zahlenrätsel: „Addiert man zum Quadrat einer Zahl das Vierfache der Zahl, erhält man $32$.“ Begründe, warum beide Lösungen passen.
\teil Rechteck: „Ein Rechteck ist $4\,$cm länger als breit und hat den Flächeninhalt $32\,\text{cm}^2$.“ Begründe, warum nur eine Lösung passt, und gib die Seitenlängen an.
\end{teile}
\end{aufgabe}
